% Options for packages loaded elsewhere
\PassOptionsToPackage{unicode}{hyperref}
\PassOptionsToPackage{hyphens}{url}
\documentclass[
]{article}
\usepackage{xcolor}
\usepackage[margin=1in]{geometry}
\usepackage{amsmath,amssymb}
\setcounter{secnumdepth}{-\maxdimen} % remove section numbering
\usepackage{iftex}
\ifPDFTeX
  \usepackage[T1]{fontenc}
  \usepackage[utf8]{inputenc}
  \usepackage{textcomp} % provide euro and other symbols
\else % if luatex or xetex
  \usepackage{unicode-math} % this also loads fontspec
  \defaultfontfeatures{Scale=MatchLowercase}
  \defaultfontfeatures[\rmfamily]{Ligatures=TeX,Scale=1}
\fi
\usepackage{lmodern}
\ifPDFTeX\else
  % xetex/luatex font selection
\fi
% Use upquote if available, for straight quotes in verbatim environments
\IfFileExists{upquote.sty}{\usepackage{upquote}}{}
\IfFileExists{microtype.sty}{% use microtype if available
  \usepackage[]{microtype}
  \UseMicrotypeSet[protrusion]{basicmath} % disable protrusion for tt fonts
}{}
\makeatletter
\@ifundefined{KOMAClassName}{% if non-KOMA class
  \IfFileExists{parskip.sty}{%
    \usepackage{parskip}
  }{% else
    \setlength{\parindent}{0pt}
    \setlength{\parskip}{6pt plus 2pt minus 1pt}}
}{% if KOMA class
  \KOMAoptions{parskip=half}}
\makeatother
\usepackage{color}
\usepackage{fancyvrb}
\newcommand{\VerbBar}{|}
\newcommand{\VERB}{\Verb[commandchars=\\\{\}]}
\DefineVerbatimEnvironment{Highlighting}{Verbatim}{commandchars=\\\{\}}
% Add ',fontsize=\small' for more characters per line
\usepackage{framed}
\definecolor{shadecolor}{RGB}{248,248,248}
\newenvironment{Shaded}{\begin{snugshade}}{\end{snugshade}}
\newcommand{\AlertTok}[1]{\textcolor[rgb]{0.94,0.16,0.16}{#1}}
\newcommand{\AnnotationTok}[1]{\textcolor[rgb]{0.56,0.35,0.01}{\textbf{\textit{#1}}}}
\newcommand{\AttributeTok}[1]{\textcolor[rgb]{0.13,0.29,0.53}{#1}}
\newcommand{\BaseNTok}[1]{\textcolor[rgb]{0.00,0.00,0.81}{#1}}
\newcommand{\BuiltInTok}[1]{#1}
\newcommand{\CharTok}[1]{\textcolor[rgb]{0.31,0.60,0.02}{#1}}
\newcommand{\CommentTok}[1]{\textcolor[rgb]{0.56,0.35,0.01}{\textit{#1}}}
\newcommand{\CommentVarTok}[1]{\textcolor[rgb]{0.56,0.35,0.01}{\textbf{\textit{#1}}}}
\newcommand{\ConstantTok}[1]{\textcolor[rgb]{0.56,0.35,0.01}{#1}}
\newcommand{\ControlFlowTok}[1]{\textcolor[rgb]{0.13,0.29,0.53}{\textbf{#1}}}
\newcommand{\DataTypeTok}[1]{\textcolor[rgb]{0.13,0.29,0.53}{#1}}
\newcommand{\DecValTok}[1]{\textcolor[rgb]{0.00,0.00,0.81}{#1}}
\newcommand{\DocumentationTok}[1]{\textcolor[rgb]{0.56,0.35,0.01}{\textbf{\textit{#1}}}}
\newcommand{\ErrorTok}[1]{\textcolor[rgb]{0.64,0.00,0.00}{\textbf{#1}}}
\newcommand{\ExtensionTok}[1]{#1}
\newcommand{\FloatTok}[1]{\textcolor[rgb]{0.00,0.00,0.81}{#1}}
\newcommand{\FunctionTok}[1]{\textcolor[rgb]{0.13,0.29,0.53}{\textbf{#1}}}
\newcommand{\ImportTok}[1]{#1}
\newcommand{\InformationTok}[1]{\textcolor[rgb]{0.56,0.35,0.01}{\textbf{\textit{#1}}}}
\newcommand{\KeywordTok}[1]{\textcolor[rgb]{0.13,0.29,0.53}{\textbf{#1}}}
\newcommand{\NormalTok}[1]{#1}
\newcommand{\OperatorTok}[1]{\textcolor[rgb]{0.81,0.36,0.00}{\textbf{#1}}}
\newcommand{\OtherTok}[1]{\textcolor[rgb]{0.56,0.35,0.01}{#1}}
\newcommand{\PreprocessorTok}[1]{\textcolor[rgb]{0.56,0.35,0.01}{\textit{#1}}}
\newcommand{\RegionMarkerTok}[1]{#1}
\newcommand{\SpecialCharTok}[1]{\textcolor[rgb]{0.81,0.36,0.00}{\textbf{#1}}}
\newcommand{\SpecialStringTok}[1]{\textcolor[rgb]{0.31,0.60,0.02}{#1}}
\newcommand{\StringTok}[1]{\textcolor[rgb]{0.31,0.60,0.02}{#1}}
\newcommand{\VariableTok}[1]{\textcolor[rgb]{0.00,0.00,0.00}{#1}}
\newcommand{\VerbatimStringTok}[1]{\textcolor[rgb]{0.31,0.60,0.02}{#1}}
\newcommand{\WarningTok}[1]{\textcolor[rgb]{0.56,0.35,0.01}{\textbf{\textit{#1}}}}
\usepackage{graphicx}
\makeatletter
\newsavebox\pandoc@box
\newcommand*\pandocbounded[1]{% scales image to fit in text height/width
  \sbox\pandoc@box{#1}%
  \Gscale@div\@tempa{\textheight}{\dimexpr\ht\pandoc@box+\dp\pandoc@box\relax}%
  \Gscale@div\@tempb{\linewidth}{\wd\pandoc@box}%
  \ifdim\@tempb\p@<\@tempa\p@\let\@tempa\@tempb\fi% select the smaller of both
  \ifdim\@tempa\p@<\p@\scalebox{\@tempa}{\usebox\pandoc@box}%
  \else\usebox{\pandoc@box}%
  \fi%
}
% Set default figure placement to htbp
\def\fps@figure{htbp}
\makeatother
\setlength{\emergencystretch}{3em} % prevent overfull lines
\providecommand{\tightlist}{%
  \setlength{\itemsep}{0pt}\setlength{\parskip}{0pt}}
\usepackage{bookmark}
\IfFileExists{xurl.sty}{\usepackage{xurl}}{} % add URL line breaks if available
\urlstyle{same}
\hypersetup{
  pdftitle={TP2},
  pdfauthor={Parfait Jemmy Prodige NGOYI, Analyste statisticien},
  hidelinks,
  pdfcreator={LaTeX via pandoc}}

\title{TP2}
\author{Parfait Jemmy Prodige NGOYI, Analyste statisticien}
\date{2026-03-17}

\begin{document}
\maketitle

\begin{enumerate}
\def\labelenumi{\arabic{enumi}.}
\setcounter{enumi}{6}
\tightlist
\item
  Chargement de sect2b\_harvestw4 et sect1\_harvestw4.
\end{enumerate}

\begin{Shaded}
\begin{Highlighting}[]
\FunctionTok{library}\NormalTok{(haven)}
\end{Highlighting}
\end{Shaded}

\begin{verbatim}
## Warning: le package 'haven' a été compilé avec la version R 4.3.3
\end{verbatim}

\begin{Shaded}
\begin{Highlighting}[]
\NormalTok{sect2\_harvest }\OtherTok{\textless{}{-}} \FunctionTok{read\_dta}\NormalTok{(}\StringTok{"sect2\_harvestw4.dta"}\NormalTok{)}
\NormalTok{sect1\_harvest }\OtherTok{\textless{}{-}} \FunctionTok{read\_dta}\NormalTok{(}\StringTok{"sect1\_harvestw4.dta"}\NormalTok{)}
\end{Highlighting}
\end{Shaded}

Jointure des bases de données

\begin{Shaded}
\begin{Highlighting}[]
\FunctionTok{library}\NormalTok{(dplyr)}
\end{Highlighting}
\end{Shaded}

\begin{verbatim}
## Warning: le package 'dplyr' a été compilé avec la version R 4.3.3
\end{verbatim}

\begin{verbatim}
## 
## Attachement du package : 'dplyr'
\end{verbatim}

\begin{verbatim}
## Les objets suivants sont masqués depuis 'package:stats':
## 
##     filter, lag
\end{verbatim}

\begin{verbatim}
## Les objets suivants sont masqués depuis 'package:base':
## 
##     intersect, setdiff, setequal, union
\end{verbatim}

\begin{Shaded}
\begin{Highlighting}[]
\NormalTok{data\_joint }\OtherTok{\textless{}{-}}\NormalTok{ sect1\_harvest }\SpecialCharTok{\%\textgreater{}\%}
  \FunctionTok{left\_join}\NormalTok{(sect2\_harvest, }\AttributeTok{by =} \FunctionTok{c}\NormalTok{(}\StringTok{"hhid"}\NormalTok{, }\StringTok{"indiv"}\NormalTok{))}
\end{Highlighting}
\end{Shaded}

Inspection les valeurs manquantes du niveau d'éducation

\begin{Shaded}
\begin{Highlighting}[]
\FunctionTok{sum}\NormalTok{(}\FunctionTok{is.na}\NormalTok{(data\_joint}\SpecialCharTok{$}\NormalTok{s2aq9))}
\end{Highlighting}
\end{Shaded}

\begin{verbatim}
## [1] 11066
\end{verbatim}

\begin{Shaded}
\begin{Highlighting}[]
\FunctionTok{mean}\NormalTok{(}\FunctionTok{is.na}\NormalTok{(data\_joint}\SpecialCharTok{$}\NormalTok{s2aq9))}
\end{Highlighting}
\end{Shaded}

\begin{verbatim}
## [1] 0.3647691
\end{verbatim}

Il y a 11066 valeurs manquantes dans la variable s2aq9 (niveau
d'éducation complété). Soit un taux de 36,48\% de valeurs manquantes sur
cette variable.

\begin{enumerate}
\def\labelenumi{\arabic{enumi}.}
\setcounter{enumi}{7}
\tightlist
\item
  Construction d'une variable ordonnée niveau\_educ à 5 catégories
  (Aucun, Primaire, Junior Secondary, Senior Secondary, Tertiaire) à
  partir de s2bq1.
\end{enumerate}

\begin{Shaded}
\begin{Highlighting}[]
\NormalTok{data\_joint }\OtherTok{\textless{}{-}}\NormalTok{ data\_joint }\SpecialCharTok{\%\textgreater{}\%}
  \FunctionTok{mutate}\NormalTok{(}\AttributeTok{niveau\_educ =} \FunctionTok{case\_when}\NormalTok{(}
\NormalTok{    s2aq9 }\SpecialCharTok{\%in\%} \FunctionTok{c}\NormalTok{(}\DecValTok{0}\NormalTok{, }\DecValTok{1}\NormalTok{, }\DecValTok{2}\NormalTok{, }\DecValTok{3}\NormalTok{, }\DecValTok{51}\NormalTok{, }\DecValTok{52}\NormalTok{, }\DecValTok{61}\NormalTok{)                    }\SpecialCharTok{\textasciitilde{}} \StringTok{"Aucun"}\NormalTok{,}
\NormalTok{    s2aq9 }\SpecialCharTok{\%in\%} \FunctionTok{c}\NormalTok{(}\DecValTok{11}\NormalTok{, }\DecValTok{12}\NormalTok{, }\DecValTok{13}\NormalTok{, }\DecValTok{14}\NormalTok{, }\DecValTok{15}\NormalTok{, }\DecValTok{16}\NormalTok{)                    }\SpecialCharTok{\textasciitilde{}} \StringTok{"Primaire"}\NormalTok{,}
\NormalTok{    s2aq9 }\SpecialCharTok{\%in\%} \FunctionTok{c}\NormalTok{(}\DecValTok{21}\NormalTok{, }\DecValTok{22}\NormalTok{, }\DecValTok{23}\NormalTok{)                                 }\SpecialCharTok{\textasciitilde{}} \StringTok{"Junior Secondary"}\NormalTok{,}
\NormalTok{    s2aq9 }\SpecialCharTok{\%in\%} \FunctionTok{c}\NormalTok{(}\DecValTok{24}\NormalTok{, }\DecValTok{25}\NormalTok{, }\DecValTok{26}\NormalTok{, }\DecValTok{27}\NormalTok{, }\DecValTok{28}\NormalTok{, }\DecValTok{31}\NormalTok{, }\DecValTok{33}\NormalTok{, }\DecValTok{34}\NormalTok{, }\DecValTok{35}\NormalTok{, }\DecValTok{321}\NormalTok{)   }\SpecialCharTok{\textasciitilde{}} \StringTok{"Senior Secondary"}\NormalTok{,}
\NormalTok{    s2aq9 }\SpecialCharTok{\%in\%} \FunctionTok{c}\NormalTok{(}\DecValTok{41}\NormalTok{, }\DecValTok{43}\NormalTok{, }\DecValTok{322}\NormalTok{, }\DecValTok{411}\NormalTok{, }\DecValTok{412}\NormalTok{, }\DecValTok{421}\NormalTok{, }\DecValTok{422}\NormalTok{, }\DecValTok{423}\NormalTok{, }\DecValTok{424}\NormalTok{)  }\SpecialCharTok{\textasciitilde{}} \StringTok{"Tertiaire"}\NormalTok{,}
    \ConstantTok{TRUE} \SpecialCharTok{\textasciitilde{}} \ConstantTok{NA\_character\_}
\NormalTok{  ),}
  \AttributeTok{niveau\_educ =} \FunctionTok{factor}\NormalTok{(niveau\_educ,}
                       \AttributeTok{levels =} \FunctionTok{c}\NormalTok{(}\StringTok{"Aucun"}\NormalTok{, }\StringTok{"Primaire"}\NormalTok{, }\StringTok{"Junior Secondary"}\NormalTok{,}
                                  \StringTok{"Senior Secondary"}\NormalTok{, }\StringTok{"Tertiaire"}\NormalTok{),}
                       \AttributeTok{ordered =} \ConstantTok{TRUE}\NormalTok{)}
\NormalTok{  )}
\end{Highlighting}
\end{Shaded}

Calcul des fréquences et proportions par catégorie

\begin{Shaded}
\begin{Highlighting}[]
\NormalTok{freq }\OtherTok{\textless{}{-}} \FunctionTok{table}\NormalTok{(data\_joint}\SpecialCharTok{$}\NormalTok{niveau\_educ)}
\NormalTok{pourcent }\OtherTok{\textless{}{-}} \FunctionTok{prop.table}\NormalTok{(freq) }\SpecialCharTok{*} \DecValTok{100}

\FunctionTok{data.frame}\NormalTok{(}
  \AttributeTok{Niveau =} \FunctionTok{names}\NormalTok{(freq),}
  \AttributeTok{Frequence =}\NormalTok{ freq,}
  \AttributeTok{Pourcentage =} \FunctionTok{round}\NormalTok{(pourcent, }\DecValTok{2}\NormalTok{)}
\NormalTok{)}
\end{Highlighting}
\end{Shaded}

\begin{verbatim}
##             Niveau   Frequence.Var1 Frequence.Freq Pourcentage.Var1
## 1            Aucun            Aucun           3368            Aucun
## 2         Primaire         Primaire           7068         Primaire
## 3 Junior Secondary Junior Secondary           2129 Junior Secondary
## 4 Senior Secondary Senior Secondary           5116 Senior Secondary
## 5        Tertiaire        Tertiaire           1590        Tertiaire
##   Pourcentage.Freq
## 1            17.48
## 2            36.68
## 3            11.05
## 4            26.55
## 5             8.25
\end{verbatim}

Visualisation avec un barplot horizontal coloré par catégorie

\begin{Shaded}
\begin{Highlighting}[]
\FunctionTok{library}\NormalTok{ (ggplot2)}
\CommentTok{\# Préparer les données avec pourcentages}
\NormalTok{data\_plot }\OtherTok{\textless{}{-}}\NormalTok{ data\_joint }\SpecialCharTok{\%\textgreater{}\%}
  \FunctionTok{filter}\NormalTok{(}\SpecialCharTok{!}\FunctionTok{is.na}\NormalTok{(niveau\_educ)) }\SpecialCharTok{\%\textgreater{}\%}
  \FunctionTok{count}\NormalTok{(niveau\_educ) }\SpecialCharTok{\%\textgreater{}\%}
  \FunctionTok{mutate}\NormalTok{(}\AttributeTok{pourcentage =}\NormalTok{ n }\SpecialCharTok{/} \FunctionTok{sum}\NormalTok{(n) }\SpecialCharTok{*} \DecValTok{100}\NormalTok{)}

\CommentTok{\# Barplot horizontal stylisé avec étiquettes}
\FunctionTok{ggplot}\NormalTok{(data\_plot, }\FunctionTok{aes}\NormalTok{(}\AttributeTok{x =}\NormalTok{ niveau\_educ, }\AttributeTok{y =}\NormalTok{ pourcentage, }\AttributeTok{fill =}\NormalTok{ niveau\_educ)) }\SpecialCharTok{+}
  \FunctionTok{geom\_bar}\NormalTok{(}\AttributeTok{stat =} \StringTok{"identity"}\NormalTok{, }\AttributeTok{width =} \FloatTok{0.7}\NormalTok{) }\SpecialCharTok{+}
  \FunctionTok{geom\_text}\NormalTok{(}\FunctionTok{aes}\NormalTok{(}\AttributeTok{label =} \FunctionTok{paste0}\NormalTok{(}\FunctionTok{round}\NormalTok{(pourcentage,}\DecValTok{1}\NormalTok{), }\StringTok{"\%"}\NormalTok{)),}
            \AttributeTok{hjust =} \SpecialCharTok{{-}}\FloatTok{0.1}\NormalTok{, }\AttributeTok{size =} \DecValTok{4}\NormalTok{) }\SpecialCharTok{+}  \CommentTok{\# étiquette à droite de la barre}
  \FunctionTok{coord\_flip}\NormalTok{() }\SpecialCharTok{+}
  \FunctionTok{scale\_fill\_brewer}\NormalTok{(}\AttributeTok{palette =} \StringTok{"Set2"}\NormalTok{) }\SpecialCharTok{+}  \CommentTok{\# palette harmonieuse}
  \FunctionTok{labs}\NormalTok{(}
    \AttributeTok{title =} \StringTok{"Distribution du niveau d\textquotesingle{}éducation"}\NormalTok{,}
    \AttributeTok{x =} \StringTok{"Niveau d\textquotesingle{}éducation"}\NormalTok{,}
    \AttributeTok{y =} \StringTok{"Pourcentage"}
\NormalTok{  ) }\SpecialCharTok{+}
  \FunctionTok{theme\_minimal}\NormalTok{(}\AttributeTok{base\_size =} \DecValTok{14}\NormalTok{) }\SpecialCharTok{+}
  \FunctionTok{theme}\NormalTok{(}
    \AttributeTok{legend.position =} \StringTok{"none"}\NormalTok{,}
    \AttributeTok{plot.title =} \FunctionTok{element\_text}\NormalTok{(}\AttributeTok{face =} \StringTok{"bold"}\NormalTok{, }\AttributeTok{size =} \DecValTok{16}\NormalTok{, }\AttributeTok{hjust =} \FloatTok{0.5}\NormalTok{),}
    \AttributeTok{axis.title.y =} \FunctionTok{element\_text}\NormalTok{(}\AttributeTok{face =} \StringTok{"bold"}\NormalTok{),}
    \AttributeTok{axis.title.x =} \FunctionTok{element\_text}\NormalTok{(}\AttributeTok{face =} \StringTok{"bold"}\NormalTok{)}
\NormalTok{  ) }\SpecialCharTok{+}
  \FunctionTok{scale\_y\_continuous}\NormalTok{(}\AttributeTok{expand =} \FunctionTok{expansion}\NormalTok{(}\AttributeTok{mult =} \FunctionTok{c}\NormalTok{(}\DecValTok{0}\NormalTok{, }\FloatTok{0.1}\NormalTok{)))}
\end{Highlighting}
\end{Shaded}

\pandocbounded{\includegraphics[keepaspectratio]{TP2_files/figure-latex/unnamed-chunk-8-1.pdf}}

La variable niveau\_educ, construite à partir de s2aq9 (sect2a), ne
couvre que les individus ayant définitivement quitté l'école. Les
individus encore scolarisés (s2bq1, sect2b) étant exclus, les résultats
sont à interpréter comme caractérisant la sous-population des sortants
définitifs, et non la population générale.

Le niveau Primaire est le plus fréquent (36,7\%), suivi du Senior
Secondary (26,5\%). Près d'un individu sur cinq (17,5\%) n'a aucun
niveau d'éducation formelle. Les niveaux Junior Secondary (11\%) et
Tertiaire (8,3\%) sont nettement moins représentés, ce qui traduit une
distribution asymétrique vers les niveaux bas. La médiane se situe
vraisemblablement au niveau Primaire.

Le faible effectif au Tertiaire est cohérent avec la construction de la
variable : les étudiants encore en cours figurent dans sect2b et sont
absents ici, ce qui sous-estime mécaniquement ce niveau.

\begin{enumerate}
\def\labelenumi{\arabic{enumi}.}
\setcounter{enumi}{8}
\tightlist
\item
  Comparaison entre la distribution du niveau d'éducation entre hommes
  et femmes (adultes 18+):
\end{enumerate}

\begin{Shaded}
\begin{Highlighting}[]
\CommentTok{\# Préparer les données correctement}
\NormalTok{profil\_educ }\OtherTok{\textless{}{-}}\NormalTok{ data\_joint }\SpecialCharTok{\%\textgreater{}\%}
  \FunctionTok{filter}\NormalTok{(}\SpecialCharTok{!}\FunctionTok{is.na}\NormalTok{(niveau\_educ), }\SpecialCharTok{!}\FunctionTok{is.na}\NormalTok{(s1q2), s1q4 }\SpecialCharTok{\textgreater{}=} \DecValTok{18}\NormalTok{) }\SpecialCharTok{\%\textgreater{}\%}
  \FunctionTok{mutate}\NormalTok{(}\AttributeTok{sexe =} \FunctionTok{factor}\NormalTok{(s1q2, }\AttributeTok{levels =} \FunctionTok{c}\NormalTok{(}\DecValTok{1}\NormalTok{,}\DecValTok{2}\NormalTok{), }\AttributeTok{labels =} \FunctionTok{c}\NormalTok{(}\StringTok{"Homme"}\NormalTok{,}\StringTok{"Femme"}\NormalTok{))) }\SpecialCharTok{\%\textgreater{}\%}
  \FunctionTok{group\_by}\NormalTok{(niveau\_educ, sexe) }\SpecialCharTok{\%\textgreater{}\%}
  \FunctionTok{summarise}\NormalTok{(}\AttributeTok{Frequence =} \FunctionTok{n}\NormalTok{(), }\AttributeTok{.groups =} \StringTok{"drop"}\NormalTok{) }\SpecialCharTok{\%\textgreater{}\%}
  \FunctionTok{group\_by}\NormalTok{(niveau\_educ) }\SpecialCharTok{\%\textgreater{}\%}
  \FunctionTok{mutate}\NormalTok{(}\AttributeTok{Pourcentage =} \FunctionTok{round}\NormalTok{(Frequence }\SpecialCharTok{/} \FunctionTok{sum}\NormalTok{(Frequence) }\SpecialCharTok{*} \DecValTok{100}\NormalTok{, }\DecValTok{2}\NormalTok{))}
\end{Highlighting}
\end{Shaded}

\begin{Shaded}
\begin{Highlighting}[]
\FunctionTok{ggplot}\NormalTok{(profil\_educ, }\FunctionTok{aes}\NormalTok{(}\AttributeTok{x =}\NormalTok{ niveau\_educ, }\AttributeTok{y =}\NormalTok{ Pourcentage, }\AttributeTok{fill =}\NormalTok{ sexe)) }\SpecialCharTok{+}
  \FunctionTok{geom\_bar}\NormalTok{(}\AttributeTok{stat =} \StringTok{"identity"}\NormalTok{) }\SpecialCharTok{+}
  \FunctionTok{geom\_text}\NormalTok{(}\FunctionTok{aes}\NormalTok{(}\AttributeTok{label =} \FunctionTok{paste0}\NormalTok{(Pourcentage, }\StringTok{"\%"}\NormalTok{)),}
            \AttributeTok{position =} \FunctionTok{position\_stack}\NormalTok{(}\AttributeTok{vjust =} \FloatTok{0.5}\NormalTok{), }\AttributeTok{size =} \DecValTok{4}\NormalTok{, }\AttributeTok{color =} \StringTok{"white"}\NormalTok{) }\SpecialCharTok{+}
  \FunctionTok{coord\_flip}\NormalTok{() }\SpecialCharTok{+}
  \FunctionTok{scale\_fill\_manual}\NormalTok{(}\AttributeTok{values =} \FunctionTok{c}\NormalTok{(}\StringTok{"steelblue"}\NormalTok{,}\StringTok{"tomato"}\NormalTok{)) }\SpecialCharTok{+}
  \FunctionTok{labs}\NormalTok{(}
    \AttributeTok{title =} \StringTok{"Répartition hommes/femmes par niveau d\textquotesingle{}éducation (18+)"}\NormalTok{,}
    \AttributeTok{x =} \StringTok{"Niveau d\textquotesingle{}éducation"}\NormalTok{,}
    \AttributeTok{y =} \StringTok{"Pourcentage"}\NormalTok{,}
    \AttributeTok{fill =} \StringTok{"Sexe"}
\NormalTok{  ) }\SpecialCharTok{+}
  \FunctionTok{theme\_minimal}\NormalTok{(}\AttributeTok{base\_size =} \DecValTok{14}\NormalTok{) }\SpecialCharTok{+}
  \FunctionTok{theme}\NormalTok{(}
    \AttributeTok{plot.title =} \FunctionTok{element\_text}\NormalTok{(}\AttributeTok{face =} \StringTok{"bold"}\NormalTok{, }\AttributeTok{hjust =} \FloatTok{0.5}\NormalTok{),}
    \AttributeTok{axis.title =} \FunctionTok{element\_text}\NormalTok{(}\AttributeTok{face =} \StringTok{"bold"}\NormalTok{)}
\NormalTok{  ) }\SpecialCharTok{+}
  \FunctionTok{scale\_y\_continuous}\NormalTok{(}\AttributeTok{expand =} \FunctionTok{expansion}\NormalTok{(}\AttributeTok{mult =} \FunctionTok{c}\NormalTok{(}\DecValTok{0}\NormalTok{,}\FloatTok{0.05}\NormalTok{)))}
\end{Highlighting}
\end{Shaded}

\pandocbounded{\includegraphics[keepaspectratio]{TP2_files/figure-latex/unnamed-chunk-10-1.pdf}}
Commentaire : On observe un gradient de genre monotone : la proportion
de femmes décroît strictement à mesure que le niveau d'éducation
augmente (53,4\% au niveau Aucun → 38,5\% au Tertiaire). Le point de
quasi-parité se situe au Junior Secondary (51,4\% F / 48,6\% H), au-delà
duquel les hommes deviennent majoritaires.

\begin{Shaded}
\begin{Highlighting}[]
\FunctionTok{library}\NormalTok{(tibble)}
\end{Highlighting}
\end{Shaded}

\begin{verbatim}
## Warning: le package 'tibble' a été compilé avec la version R 4.3.3
\end{verbatim}

\begin{Shaded}
\begin{Highlighting}[]
\CommentTok{\# Filtrer adultes 18+ et supprimer les NA}
\NormalTok{table\_edu\_sexe }\OtherTok{\textless{}{-}}\NormalTok{ data\_joint }\SpecialCharTok{\%\textgreater{}\%}
  \FunctionTok{filter}\NormalTok{(}\SpecialCharTok{!}\FunctionTok{is.na}\NormalTok{(niveau\_educ), }\SpecialCharTok{!}\FunctionTok{is.na}\NormalTok{(s1q2), s1q4 }\SpecialCharTok{\textgreater{}=} \DecValTok{18}\NormalTok{) }\SpecialCharTok{\%\textgreater{}\%}
  \FunctionTok{mutate}\NormalTok{(}\AttributeTok{sexe =} \FunctionTok{factor}\NormalTok{(s1q2, }\AttributeTok{levels =} \FunctionTok{c}\NormalTok{(}\DecValTok{1}\NormalTok{,}\DecValTok{2}\NormalTok{), }\AttributeTok{labels =} \FunctionTok{c}\NormalTok{(}\StringTok{"Homme"}\NormalTok{,}\StringTok{"Femme"}\NormalTok{))) }\SpecialCharTok{\%\textgreater{}\%}
  \FunctionTok{count}\NormalTok{(niveau\_educ, sexe) }\SpecialCharTok{\%\textgreater{}\%}
\NormalTok{  tidyr}\SpecialCharTok{::}\FunctionTok{pivot\_wider}\NormalTok{(}\AttributeTok{names\_from =}\NormalTok{ sexe, }\AttributeTok{values\_from =}\NormalTok{ n, }\AttributeTok{values\_fill =} \DecValTok{0}\NormalTok{) }\SpecialCharTok{\%\textgreater{}\%}
  \FunctionTok{column\_to\_rownames}\NormalTok{(}\StringTok{"niveau\_educ"}\NormalTok{)}

\CommentTok{\# Transformer en tableau classique}
\NormalTok{table\_chi }\OtherTok{\textless{}{-}} \FunctionTok{as.matrix}\NormalTok{(table\_edu\_sexe)}
\NormalTok{table\_chi}
\end{Highlighting}
\end{Shaded}

\begin{verbatim}
##                  Homme Femme
## Aucun              502   574
## Primaire          1198  1416
## Junior Secondary   359   379
## Senior Secondary  2402  1939
## Tertiaire          970   606
\end{verbatim}

\begin{Shaded}
\begin{Highlighting}[]
\CommentTok{\# Test de Chi{-}2}
\NormalTok{chi\_test }\OtherTok{\textless{}{-}} \FunctionTok{chisq.test}\NormalTok{(table\_chi)}
\NormalTok{chi\_test}
\end{Highlighting}
\end{Shaded}

\begin{verbatim}
## 
##  Pearson's Chi-squared test
## 
## data:  table_chi
## X-squared = 131.48, df = 4, p-value < 2.2e-16
\end{verbatim}

Le résultat est hautement significatif (p \textless{} 0,001) --- on
rejette sans ambiguïté l'hypothèse nulle d'indépendance entre le sexe et
le niveau d'éducation. Il existe une association statistiquement
significative entre le niveau d'éducation et le sexe.

\begin{Shaded}
\begin{Highlighting}[]
\NormalTok{data\_plot }\OtherTok{\textless{}{-}}\NormalTok{ data\_joint }\SpecialCharTok{\%\textgreater{}\%}
  \FunctionTok{filter}\NormalTok{(}\SpecialCharTok{!}\FunctionTok{is.na}\NormalTok{(niveau\_educ), }\SpecialCharTok{!}\FunctionTok{is.na}\NormalTok{(s1q4))}

\FunctionTok{ggplot}\NormalTok{(data\_plot, }\FunctionTok{aes}\NormalTok{(}\AttributeTok{x =}\NormalTok{ niveau\_educ, }\AttributeTok{y =}\NormalTok{ s1q4, }\AttributeTok{fill =}\NormalTok{ niveau\_educ)) }\SpecialCharTok{+}
  \FunctionTok{geom\_boxplot}\NormalTok{(}\AttributeTok{outlier.shape =} \DecValTok{21}\NormalTok{, }\AttributeTok{outlier.size =} \DecValTok{2}\NormalTok{, }\AttributeTok{alpha =} \FloatTok{0.8}\NormalTok{, }\AttributeTok{width =} \FloatTok{0.5}\NormalTok{) }\SpecialCharTok{+}
  \FunctionTok{scale\_fill\_brewer}\NormalTok{(}\AttributeTok{palette =} \StringTok{"Set2"}\NormalTok{) }\SpecialCharTok{+}
  \FunctionTok{labs}\NormalTok{(}
    \AttributeTok{title =} \StringTok{"Distribution de l\textquotesingle{}âge par niveau d\textquotesingle{}éducation"}\NormalTok{,}
    \AttributeTok{x =} \StringTok{"Niveau d\textquotesingle{}éducation"}\NormalTok{,}
    \AttributeTok{y =} \StringTok{"Âge"}
\NormalTok{  ) }\SpecialCharTok{+}
  \FunctionTok{theme\_minimal}\NormalTok{(}\AttributeTok{base\_size =} \DecValTok{13}\NormalTok{) }\SpecialCharTok{+}
  \FunctionTok{theme}\NormalTok{(}
    \AttributeTok{legend.position =} \StringTok{"none"}\NormalTok{,}
    \AttributeTok{plot.title =} \FunctionTok{element\_text}\NormalTok{(}\AttributeTok{face =} \StringTok{"bold"}\NormalTok{, }\AttributeTok{hjust =} \FloatTok{0.5}\NormalTok{),}
    \AttributeTok{panel.grid.minor =} \FunctionTok{element\_blank}\NormalTok{()}
\NormalTok{  )}
\end{Highlighting}
\end{Shaded}

\pandocbounded{\includegraphics[keepaspectratio]{TP2_files/figure-latex/unnamed-chunk-13-1.pdf}}
Une progression cohérente mais non monotone. Les médianes croissent
globalement avec le niveau d'éducation, ce qui est attendu --- les
individus plus éduqués sont structurellement plus âgés au moment de la
sortie définitive.

Le cas ``Aucun'' est atypique. La médiane très basse (\textasciitilde8
ans) et la très grande dispersion (IQR large, valeurs jusqu'à 100 ans)
traduisent une population hétérogène : d'un côté des enfants n'ayant
jamais été scolarisés, de l'autre des personnes âgées dont la
non-scolarisation reflète des générations sans accès à l'école. C'est la
boîte la plus dispersée de toutes.

Le Junior Secondary se distingue par une boîte très étroite et une
médiane basse (\textasciitilde17 ans) --- c'est le niveau le plus
homogène en âge, ce qui suggère que les sortants à ce niveau
correspondent à une tranche d'âge précise et concentrée.

Le Tertiaire présente la médiane la plus élevée (\textasciitilde33 ans)
avec une dispersion modérée, cohérent avec une sortie tardive après de
longues études --- mais là encore, les étudiants encore en cours sont
absents, ce qui tire la médiane vers le haut (seuls les plus anciens
diplômés sont captés). Les valeurs aberrantes (points au-delà des
moustaches) sont présentes dans tous les groupes, particulièrement pour
Aucun et Primaire, avec des individus âgés de 75 à 100 ans --- des
générations pour lesquelles la scolarisation était rare ou inexistante.

Kruskal-Wallis + post-hoc Dunn.

\begin{Shaded}
\begin{Highlighting}[]
\FunctionTok{install.packages}\NormalTok{(}\StringTok{"dunn.test"}\NormalTok{, }\AttributeTok{repos =} \StringTok{"https://cloud.r{-}project.org"}\NormalTok{)}
\end{Highlighting}
\end{Shaded}

\begin{verbatim}
## Installation du package dans 'C:/Users/LENOVO T470S/AppData/Local/R/win-library/4.3'
## (car 'lib' n'est pas spécifié)
\end{verbatim}

\begin{verbatim}
## 
##   Une version binaire est disponible mais la version du source est plus
##   récente:
##           binary source needs_compilation
## dunn.test  1.3.6  1.3.7             FALSE
\end{verbatim}

\begin{verbatim}
## installation du package source 'dunn.test'
\end{verbatim}

\begin{Shaded}
\begin{Highlighting}[]
\FunctionTok{library}\NormalTok{(dunn.test)}
\end{Highlighting}
\end{Shaded}

\begin{Shaded}
\begin{Highlighting}[]
\CommentTok{\# Kruskal{-}Wallis}
\FunctionTok{kruskal.test}\NormalTok{(s1q4 }\SpecialCharTok{\textasciitilde{}}\NormalTok{ niveau\_educ, }\AttributeTok{data =}\NormalTok{ data\_plot)}
\end{Highlighting}
\end{Shaded}

\begin{verbatim}
## 
##  Kruskal-Wallis rank sum test
## 
## data:  s1q4 by niveau_educ
## Kruskal-Wallis chi-squared = 3928.3, df = 4, p-value < 2.2e-16
\end{verbatim}

\begin{Shaded}
\begin{Highlighting}[]
\CommentTok{\# Post{-}hoc Dunn avec correction de Bonferroni}
\FunctionTok{dunn.test}\NormalTok{(data\_plot}\SpecialCharTok{$}\NormalTok{s1q4, }
          \AttributeTok{g =}\NormalTok{ data\_plot}\SpecialCharTok{$}\NormalTok{niveau\_educ,}
          \AttributeTok{method =} \StringTok{"bonferroni"}\NormalTok{,}
          \AttributeTok{kw =} \ConstantTok{TRUE}\NormalTok{,}
          \AttributeTok{label =} \ConstantTok{TRUE}\NormalTok{)}
\end{Highlighting}
\end{Shaded}

\begin{verbatim}
##   Kruskal-Wallis rank sum test
## 
## data: x and group
## Kruskal-Wallis chi-squared = 3928.2908, df = 4, p-value = 0
## 
## 
##                    Dunn's Pairwise Comparison of x by group                   
##                                  (Bonferroni)                                 
## Col Mean-│
## Row Mean │      Aucun   Junior S   Primaire   Senior S
## ─────────┼────────────────────────────────────────────
## Junior S │  -17.99966
##          │     0.0000*
##          │
## Primaire │  -18.57679   4.425702
##          │     0.0000*    0.0000*
##          │
## Senior S │  -49.64826  -23.39207  -38.82706
##          │     0.0000*    0.0000*    0.0000*
##          │
## Tertiair │  -46.96376  -28.07685  -37.47034  -11.40012
##          │     0.0000*    0.0000*    0.0000*    0.0000*
## 
## α = 0.05
## Reject Ho if p ≤ α/2, where p = Pr(Z ≥ |z|)
\end{verbatim}

La p-valeur quasi-nulle conduit au rejet de l'hypothèse nulle d'égalité
des distributions d'âge entre les niveaux d'éducation. Il existe au
moins une paire de groupes dont les distributions d'âge diffèrent
significativement.

Toutes les paires sont significativement différentes, ce qui confirme
que chaque niveau d'éducation correspond à une distribution d'âge
distincte dans la population des sortants définitifs.

L'anomalie Aucun/Primaire mérite attention : la statistique Z de
Primaire vs Aucun (-18,58) indique que les individus du niveau Primaire
sont en médiane plus âgés que ceux du niveau Aucun --- contra-intuitif a
priori, mais cohérent avec le boxplot : le groupe ``Aucun'' contient
beaucoup d'enfants très jeunes n'ayant jamais été scolarisés, ce qui
tire la médiane vers le bas.

La comparaison Junior Secondary vs Primaire (Z = +4,43) est la plus
faible de toutes, bien que significative --- les deux groupes sont
relativement proches en âge, ce qui correspond aux boîtes adjacentes et
de taille similaire observées sur le graphique. Les écarts les plus
forts impliquent Senior Secondary et Tertiaire vs Aucun (Z ≈ -50 et
-47), traduisant des différences de rangs massives entre les extrêmes de
la distribution éducative.

\begin{enumerate}
\def\labelenumi{\arabic{enumi}.}
\setcounter{enumi}{10}
\tightlist
\item
  Comparaison du taux de scolarisation des 6-17 ans entre zones rurales
  et urbaines
\end{enumerate}

\begin{Shaded}
\begin{Highlighting}[]
\CommentTok{\# 1. Préparer les données}
\NormalTok{data\_scol }\OtherTok{\textless{}{-}}\NormalTok{ data\_joint }\SpecialCharTok{\%\textgreater{}\%}
  \FunctionTok{filter}\NormalTok{(s1q4 }\SpecialCharTok{\textgreater{}=} \DecValTok{6} \SpecialCharTok{\&}\NormalTok{ s1q4 }\SpecialCharTok{\textless{}=} \DecValTok{17}\NormalTok{, }\SpecialCharTok{!}\FunctionTok{is.na}\NormalTok{(sector.x), }\SpecialCharTok{!}\FunctionTok{is.na}\NormalTok{(s2aq6)) }\SpecialCharTok{\%\textgreater{}\%}
  \FunctionTok{mutate}\NormalTok{(}
    \AttributeTok{zone =} \FunctionTok{factor}\NormalTok{(sector.x, }\AttributeTok{levels =} \FunctionTok{c}\NormalTok{(}\DecValTok{1}\NormalTok{, }\DecValTok{2}\NormalTok{), }\AttributeTok{labels =} \FunctionTok{c}\NormalTok{(}\StringTok{"Urbain"}\NormalTok{, }\StringTok{"Rural"}\NormalTok{)),}
    \AttributeTok{scolarise =} \FunctionTok{as.numeric}\NormalTok{(s2aq6 }\SpecialCharTok{==} \DecValTok{1}\NormalTok{)}
\NormalTok{  )}

\CommentTok{\# 2. Tableau de contingence}
\NormalTok{table\_scol }\OtherTok{\textless{}{-}} \FunctionTok{table}\NormalTok{(data\_scol}\SpecialCharTok{$}\NormalTok{zone, data\_scol}\SpecialCharTok{$}\NormalTok{scolarise,}
                    \AttributeTok{dnn =} \FunctionTok{c}\NormalTok{(}\StringTok{"Zone"}\NormalTok{, }\StringTok{"Scolarisé"}\NormalTok{))}
\FunctionTok{print}\NormalTok{(table\_scol)}
\end{Highlighting}
\end{Shaded}

\begin{verbatim}
##         Scolarisé
## Zone        0    1
##   Urbain   98 2315
##   Rural  1212 5515
\end{verbatim}

\begin{Shaded}
\begin{Highlighting}[]
\CommentTok{\# 3. Test du Chi{-}deux}
\FunctionTok{chisq.test}\NormalTok{(table\_scol)}
\end{Highlighting}
\end{Shaded}

\begin{verbatim}
## 
##  Pearson's Chi-squared test with Yates' continuity correction
## 
## data:  table_scol
## X-squared = 280.57, df = 1, p-value < 2.2e-16
\end{verbatim}

\begin{Shaded}
\begin{Highlighting}[]
\CommentTok{\# 4. Taux + IC à 95\% par zone}
\NormalTok{data\_ic }\OtherTok{\textless{}{-}}\NormalTok{ data\_scol }\SpecialCharTok{\%\textgreater{}\%}
  \FunctionTok{group\_by}\NormalTok{(zone) }\SpecialCharTok{\%\textgreater{}\%}
  \FunctionTok{summarise}\NormalTok{(}
    \AttributeTok{n        =} \FunctionTok{n}\NormalTok{(),}
    \AttributeTok{n\_scol   =} \FunctionTok{sum}\NormalTok{(scolarise),}
    \AttributeTok{taux     =}\NormalTok{ n\_scol }\SpecialCharTok{/}\NormalTok{ n,}
    \AttributeTok{se       =} \FunctionTok{sqrt}\NormalTok{(taux }\SpecialCharTok{*}\NormalTok{ (}\DecValTok{1} \SpecialCharTok{{-}}\NormalTok{ taux) }\SpecialCharTok{/}\NormalTok{ n),}
    \AttributeTok{ic\_lower =}\NormalTok{ taux }\SpecialCharTok{{-}} \FloatTok{1.96} \SpecialCharTok{*}\NormalTok{ se,}
    \AttributeTok{ic\_upper =}\NormalTok{ taux }\SpecialCharTok{+} \FloatTok{1.96} \SpecialCharTok{*}\NormalTok{ se}
\NormalTok{  )}

\FunctionTok{print}\NormalTok{(data\_ic)}
\end{Highlighting}
\end{Shaded}

\begin{verbatim}
## # A tibble: 2 x 7
##   zone       n n_scol  taux      se ic_lower ic_upper
##   <fct>  <int>  <dbl> <dbl>   <dbl>    <dbl>    <dbl>
## 1 Urbain  2413   2315 0.959 0.00402    0.952    0.967
## 2 Rural   6727   5515 0.820 0.00469    0.811    0.829
\end{verbatim}

\begin{Shaded}
\begin{Highlighting}[]
\CommentTok{\# 5. Graphique en barres groupées avec IC 95\%}
\FunctionTok{ggplot}\NormalTok{(data\_ic, }\FunctionTok{aes}\NormalTok{(}\AttributeTok{x =}\NormalTok{ zone, }\AttributeTok{y =}\NormalTok{ taux }\SpecialCharTok{*} \DecValTok{100}\NormalTok{, }\AttributeTok{fill =}\NormalTok{ zone)) }\SpecialCharTok{+}
  \FunctionTok{geom\_col}\NormalTok{(}\AttributeTok{width =} \FloatTok{0.5}\NormalTok{, }\AttributeTok{alpha =} \FloatTok{0.85}\NormalTok{) }\SpecialCharTok{+}
  \FunctionTok{geom\_errorbar}\NormalTok{(}\FunctionTok{aes}\NormalTok{(}\AttributeTok{ymin =}\NormalTok{ ic\_lower }\SpecialCharTok{*} \DecValTok{100}\NormalTok{, }\AttributeTok{ymax =}\NormalTok{ ic\_upper }\SpecialCharTok{*} \DecValTok{100}\NormalTok{),}
                \AttributeTok{width =} \FloatTok{0.15}\NormalTok{, }\AttributeTok{linewidth =} \FloatTok{0.8}\NormalTok{) }\SpecialCharTok{+}
  \FunctionTok{geom\_text}\NormalTok{(}\FunctionTok{aes}\NormalTok{(}\AttributeTok{label =} \FunctionTok{paste0}\NormalTok{(}\FunctionTok{round}\NormalTok{(taux }\SpecialCharTok{*} \DecValTok{100}\NormalTok{, }\DecValTok{1}\NormalTok{), }\StringTok{"\%"}\NormalTok{)),}
            \AttributeTok{vjust =} \SpecialCharTok{{-}}\FloatTok{1.8}\NormalTok{, }\AttributeTok{fontface =} \StringTok{"bold"}\NormalTok{, }\AttributeTok{size =} \FloatTok{3.5}\NormalTok{) }\SpecialCharTok{+}
  \FunctionTok{scale\_fill\_brewer}\NormalTok{(}\AttributeTok{palette =} \StringTok{"Set2"}\NormalTok{) }\SpecialCharTok{+}
  \FunctionTok{scale\_y\_continuous}\NormalTok{(}\AttributeTok{limits =} \FunctionTok{c}\NormalTok{(}\DecValTok{0}\NormalTok{, }\DecValTok{100}\NormalTok{),}
                     \AttributeTok{labels =}\NormalTok{ scales}\SpecialCharTok{::}\FunctionTok{percent\_format}\NormalTok{(}\AttributeTok{scale =} \DecValTok{1}\NormalTok{)) }\SpecialCharTok{+}
  \FunctionTok{labs}\NormalTok{(}
    \AttributeTok{title    =} \StringTok{"Taux de scolarisation des 6{-}17 ans}\SpecialCharTok{\textbackslash{}n}\StringTok{par zone géographique"}\NormalTok{,}
    \AttributeTok{x        =} \StringTok{"Zone"}\NormalTok{,}
    \AttributeTok{y        =} \StringTok{"Taux de scolarisation (\%)"}\NormalTok{,}
    \AttributeTok{caption  =} \StringTok{"IC à 95\% — Test du Chi{-}deux"}
\NormalTok{  ) }\SpecialCharTok{+}
  \FunctionTok{theme\_minimal}\NormalTok{(}\AttributeTok{base\_size =} \DecValTok{13}\NormalTok{) }\SpecialCharTok{+}
  \FunctionTok{theme}\NormalTok{(}
    \AttributeTok{legend.position  =} \StringTok{"none"}\NormalTok{,}
    \AttributeTok{plot.title       =} \FunctionTok{element\_text}\NormalTok{(}\AttributeTok{face =} \StringTok{"bold"}\NormalTok{, }\AttributeTok{hjust =} \FloatTok{0.5}\NormalTok{),}
    \AttributeTok{panel.grid.minor =} \FunctionTok{element\_blank}\NormalTok{()}
\NormalTok{  )}
\end{Highlighting}
\end{Shaded}

\pandocbounded{\includegraphics[keepaspectratio]{TP2_files/figure-latex/unnamed-chunk-16-1.pdf}}
L'écart de scolarisation entre zones urbaines et rurales est hautement
significatif (p \textless{} 0,001) et substantiel : les enfants de 6-17
ans en zone urbaine ont un taux de scolarisation de 95,9\%, contre
82,0\% en zone rurale, soit un écart de près de 14 points de
pourcentage. Les intervalles de confiance ne se chevauchent pas, ce qui
confirme visuellement la significativité de l'écart indépendamment du
test formel.

Par aileurs, le taux de scolarisation des 6-17 ans est plus élévé en
milieu urbain (95,9\%) qu'en milieu rural (82\%).

\begin{enumerate}
\def\labelenumi{\arabic{enumi}.}
\setcounter{enumi}{9}
\tightlist
\item
  Heatmap (geom\_tile) croisant État nigérian (state) et niveau
  d'éducation: part d'adultes sans instruction par État, colorée par
  quintile.
\end{enumerate}

\begin{Shaded}
\begin{Highlighting}[]
\NormalTok{data\_joint }\OtherTok{\textless{}{-}}\NormalTok{ data\_joint }\SpecialCharTok{\%\textgreater{}\%}
  \FunctionTok{mutate}\NormalTok{(}\AttributeTok{state\_nom =} \FunctionTok{factor}\NormalTok{(state.x,}
    \AttributeTok{levels =} \DecValTok{1}\SpecialCharTok{:}\DecValTok{37}\NormalTok{,}
    \AttributeTok{labels =} \FunctionTok{c}\NormalTok{(}\StringTok{"Abia"}\NormalTok{, }\StringTok{"Adamawa"}\NormalTok{, }\StringTok{"Akwa Ibom"}\NormalTok{, }\StringTok{"Anambra"}\NormalTok{, }\StringTok{"Bauchi"}\NormalTok{,}
               \StringTok{"Bayelsa"}\NormalTok{, }\StringTok{"Benue"}\NormalTok{, }\StringTok{"Borno"}\NormalTok{, }\StringTok{"Cross River"}\NormalTok{, }\StringTok{"Delta"}\NormalTok{,}
               \StringTok{"Ebonyi"}\NormalTok{, }\StringTok{"Edo"}\NormalTok{, }\StringTok{"Ekiti"}\NormalTok{, }\StringTok{"Enugu"}\NormalTok{, }\StringTok{"Gombe"}\NormalTok{,}
               \StringTok{"Imo"}\NormalTok{, }\StringTok{"Jigawa"}\NormalTok{, }\StringTok{"Kaduna"}\NormalTok{, }\StringTok{"Kano"}\NormalTok{, }\StringTok{"Katsina"}\NormalTok{,}
               \StringTok{"Kebbi"}\NormalTok{, }\StringTok{"Kogi"}\NormalTok{, }\StringTok{"Kwara"}\NormalTok{, }\StringTok{"Lagos"}\NormalTok{, }\StringTok{"Nasarawa"}\NormalTok{,}
               \StringTok{"Niger"}\NormalTok{, }\StringTok{"Ogun"}\NormalTok{, }\StringTok{"Ondo"}\NormalTok{, }\StringTok{"Osun"}\NormalTok{, }\StringTok{"Oyo"}\NormalTok{,}
               \StringTok{"Plateau"}\NormalTok{, }\StringTok{"Rivers"}\NormalTok{, }\StringTok{"Sokoto"}\NormalTok{, }\StringTok{"Taraba"}\NormalTok{, }\StringTok{"Yobe"}\NormalTok{,}
               \StringTok{"Zamfara"}\NormalTok{, }\StringTok{"FCT"}\NormalTok{)}
\NormalTok{  ))}
\end{Highlighting}
\end{Shaded}

\begin{Shaded}
\begin{Highlighting}[]
\NormalTok{data\_heat }\OtherTok{\textless{}{-}}\NormalTok{ data\_joint }\SpecialCharTok{\%\textgreater{}\%}
  \FunctionTok{filter}\NormalTok{(s1q4 }\SpecialCharTok{\textgreater{}=} \DecValTok{18}\NormalTok{, }\SpecialCharTok{!}\FunctionTok{is.na}\NormalTok{(state\_nom), }\SpecialCharTok{!}\FunctionTok{is.na}\NormalTok{(niveau\_educ)) }\SpecialCharTok{\%\textgreater{}\%}
  \FunctionTok{group\_by}\NormalTok{(state\_nom) }\SpecialCharTok{\%\textgreater{}\%}
  \FunctionTok{summarise}\NormalTok{(}
    \AttributeTok{n          =} \FunctionTok{n}\NormalTok{(),}
    \AttributeTok{n\_aucun    =} \FunctionTok{sum}\NormalTok{(niveau\_educ }\SpecialCharTok{==} \StringTok{"Aucun"}\NormalTok{),}
    \AttributeTok{part\_aucun =}\NormalTok{ n\_aucun }\SpecialCharTok{/}\NormalTok{ n}
\NormalTok{  ) }\SpecialCharTok{\%\textgreater{}\%}
  \FunctionTok{mutate}\NormalTok{(}
    \AttributeTok{quintile =} \FunctionTok{cut}\NormalTok{(part\_aucun,}
                   \AttributeTok{breaks =} \FunctionTok{quantile}\NormalTok{(part\_aucun, }\AttributeTok{probs =} \FunctionTok{seq}\NormalTok{(}\DecValTok{0}\NormalTok{, }\DecValTok{1}\NormalTok{, }\FloatTok{0.2}\NormalTok{)),}
                   \AttributeTok{labels =} \FunctionTok{c}\NormalTok{(}\StringTok{"Q1 (plus bas)"}\NormalTok{, }\StringTok{"Q2"}\NormalTok{, }\StringTok{"Q3"}\NormalTok{, }\StringTok{"Q4"}\NormalTok{, }\StringTok{"Q5 (plus haut)"}\NormalTok{),}
                   \AttributeTok{include.lowest =} \ConstantTok{TRUE}\NormalTok{)}
\NormalTok{  )}

\FunctionTok{ggplot}\NormalTok{(data\_heat, }\FunctionTok{aes}\NormalTok{(}\AttributeTok{x =} \DecValTok{1}\NormalTok{, }\AttributeTok{y =} \FunctionTok{reorder}\NormalTok{(state\_nom, part\_aucun), }\AttributeTok{fill =}\NormalTok{ quintile)) }\SpecialCharTok{+}
  \FunctionTok{geom\_tile}\NormalTok{(}\AttributeTok{color =} \StringTok{"white"}\NormalTok{, }\AttributeTok{linewidth =} \FloatTok{0.5}\NormalTok{) }\SpecialCharTok{+}
  \FunctionTok{geom\_text}\NormalTok{(}\FunctionTok{aes}\NormalTok{(}\AttributeTok{label =} \FunctionTok{paste0}\NormalTok{(}\FunctionTok{round}\NormalTok{(part\_aucun }\SpecialCharTok{*} \DecValTok{100}\NormalTok{, }\DecValTok{1}\NormalTok{), }\StringTok{"\%"}\NormalTok{)),}
            \AttributeTok{size =} \DecValTok{3}\NormalTok{, }\AttributeTok{color =} \StringTok{"white"}\NormalTok{, }\AttributeTok{fontface =} \StringTok{"bold"}\NormalTok{) }\SpecialCharTok{+}
  \FunctionTok{scale\_fill\_brewer}\NormalTok{(}\AttributeTok{palette =} \StringTok{"YlOrRd"}\NormalTok{, }\AttributeTok{direction =} \DecValTok{1}\NormalTok{) }\SpecialCharTok{+}
  \FunctionTok{labs}\NormalTok{(}
    \AttributeTok{title    =} \StringTok{"Part d\textquotesingle{}adultes sans instruction par État nigérian"}\NormalTok{,}
    \AttributeTok{subtitle =} \StringTok{"Colorée par quintile — adultes 18 ans et plus"}\NormalTok{,}
    \AttributeTok{x        =} \ConstantTok{NULL}\NormalTok{,}
    \AttributeTok{y        =} \StringTok{"État"}\NormalTok{,}
    \AttributeTok{fill     =} \StringTok{"Quintile"}
\NormalTok{  ) }\SpecialCharTok{+}
  \FunctionTok{theme\_minimal}\NormalTok{(}\AttributeTok{base\_size =} \DecValTok{12}\NormalTok{) }\SpecialCharTok{+}
  \FunctionTok{theme}\NormalTok{(}
    \AttributeTok{axis.text.x     =} \FunctionTok{element\_blank}\NormalTok{(),}
    \AttributeTok{axis.ticks.x    =} \FunctionTok{element\_blank}\NormalTok{(),}
    \AttributeTok{panel.grid      =} \FunctionTok{element\_blank}\NormalTok{(),}
    \AttributeTok{plot.title      =} \FunctionTok{element\_text}\NormalTok{(}\AttributeTok{face =} \StringTok{"bold"}\NormalTok{, }\AttributeTok{hjust =} \FloatTok{0.5}\NormalTok{),}
    \AttributeTok{plot.subtitle   =} \FunctionTok{element\_text}\NormalTok{(}\AttributeTok{hjust =} \FloatTok{0.5}\NormalTok{, }\AttributeTok{color =} \StringTok{"grey40"}\NormalTok{),}
    \AttributeTok{legend.position =} \StringTok{"right"}
\NormalTok{  )}
\end{Highlighting}
\end{Shaded}

\pandocbounded{\includegraphics[keepaspectratio]{TP2_files/figure-latex/unnamed-chunk-18-1.pdf}}

\end{document}
